\section{Finite, countable, and uncountable sets}
\label{sec:set-prerequisites-05}

\begin{definition}[Finite set]
A set \(A\) is \emph{finite} iff there exists \(n\in\N\) and a bijection from
\(n=\set{0,\ldots,n-1}\) to \(A\).
\end{definition}

\begin{definition}[Countable set]
A set \(A\) is \emph{countable} iff there exists an injection \(A\to\N\).
It is \emph{countably infinite} iff there exists a bijection \(\N\to A\).
\end{definition}

This document includes finite sets among countable sets.  The convention is
used in the countable-language versions of the completeness and omitting-types
theorems.

\begin{definition}[Cardinal comparison]
For sets \(A\) and \(B\), write \(\abs A\leq\abs B\) iff there exists an
injection \(A\to B\), and write \(\abs A=\abs B\) iff there exists a bijection
\(A\to B\).  Write \(\abs A<\abs B\) iff
\(\abs A\leq\abs B\) and \(\abs A\ne\abs B\).
\end{definition}

\begin{theorem}[Cantor--Schr\"oder--Bernstein]
If there is an injection \(f:A\to B\) and an injection \(g:B\to A\), then
there is a bijection \(h:A\to B\).
\end{theorem}
\begin{proof}
Let \(A_0=A\setminus g[B]\), and recursively let
\(A_{n+1}=g[f[A_n]]\).  Put \(C=\bigcup_{n\in\N}A_n\).  Define
\[
  h(a)=
  \begin{cases}
    f(a),&a\in C,\\
    g^{-1}(a),&a\notin C.
  \end{cases}
\]
The second case is defined because \(A_0\subseteq C\), so
\(a\notin C\) implies \(a\in g[B]\), and injectivity of \(g\) makes the
preimage unique.  The two ranges are disjoint: if \(a\in C\), then
\(g(f(a))\in C\); if \(a'\notin C\), then \(a'=g(b)\) with
\(g(b)\notin C\), so \(b\notin f[C]\).  Each branch is injective.  For
surjectivity, let \(b\in B\).  If \(b\in f[C]\), it is in the first range.
Otherwise \(g(b)\notin C\): since \(g(b)\in g[B]\), the element
\(g(b)\) is not in \(A_0=A\setminus g[B]\); and if
\(g(b)\in A_{n+1}=g[f[A_n]]\), injectivity of \(g\) would give
\(b\in f[C]\), contrary to the case assumption.  Thus
\(h(g(b))=b\).  Hence \(h\) is bijective.
\end{proof}

\begin{theorem}[Cantor]
For every set \(A\), there is no surjection \(f:A\to\Pow(A)\).  Consequently,
\(\abs A<\abs{\Pow(A)}\).
\end{theorem}
\begin{proof}
Suppose \(f:A\to\Pow(A)\).  Define
\[
  D\defeq\set{a\in A:a\notin f(a)}.
\]
If \(f\) were surjective, there would be \(d\in A\) with \(f(d)=D\).  Then
\[
  d\in D\quad\Longleftrightarrow\quad d\notin f(d)
  \quad\Longleftrightarrow\quad d\notin D,
\]
a contradiction.  Thus no such surjection exists.  The singleton map
\(a\mapsto\set a\) is an injection \(A\to\Pow(A)\), so the cardinal inequality
is strict.
\end{proof}

